\documentclass[../../../include/open-logic-section]{subfiles}

\begin{document}

\olfileid[hi]{sth}{cardinals}{classing}
\olsection{परिमित, \usetoken{S}{enumerable}, \usetoken{S}{nonenumerable}}

अब जबकि कार्डिनलों से हमारा परिचय हो गया है, कार्डिनलों की विभिन्न
किस्मों पर कुछ समय चर्चा करना उपयोगी है; विशेषतः परिमित, !!{enumerable}
और !!{nonenumerable} कार्डिनलों पर।

हमारे पहले दो परिणामों से निष्पन्न होगा कि परिमित कार्डिनल ठीक परिमित
क्रमसूचक होंगे, जिन्हें हमने \olref[z][infinity-again]{defnomega} में
अपनी \emph{प्राकृतिक संख्याओं} के रूप में परिभाषित किया था:

\begin{prop}\ollabel{finitecardisoequal}
मान लें $n,m\in\omega$। तब $n=m$ तभी जब $\cardeq{n}{m}$।
\end{prop}

\begin{proof}
\emph{बाएँ-से-दाएँ} दिशा तुच्छ है। \emph{दाएँ-से-बाएँ} दिशा सिद्ध करने के
लिए मान लें कि $n\neq m$ होते हुए भी $\cardeq{n}{m}$। त्रिविभाजन से या तो
$n\in m$ या $m\in n$; व्यापकता खोए बिना $n\in m$ मान लें। तब
$n\subsetneq m$ और कोई !!a{bijection} $f\colon m\to n$ है, इसलिए $m$
डेडेकिंड-अनंत है, जो
\olref[z][infinity-again]{naturalnumbersarentinfinite} के विरुद्ध है।
\end{proof}

\begin{cor}\ollabel{naturalsarecardinals}
यदि $n\in\omega$, तो $n$ एक कार्डिनल है।
\end{cor}

\begin{proof}
तत्काल।
\end{proof}
\noindent
इससे यह भी निष्पन्न होता है कि किसी कार्डिनल को ``परिमित'' या ``अनंत'' कहने
के अर्थ की कई उचित धारणाएँ परस्पर समान हैं:
\begin{thm}\ollabel{generalinfinitycharacter}किसी भी समुच्चय $A$ के लिए निम्न कथन तुल्य हैं:
\begin{enumerate}
	\item\ollabel{card:notinomega} $\card{A} \notin \omega$, अर्थात्
	$A$ कोई प्राकृतिक संख्या नहीं है;
	\item\ollabel{card:omegaplus} $\omega \leq \card{A}$;
	\item\ollabel{card:infinite} $A$ डेडेकिंड-अनंत है।
\end{enumerate}
\end{thm}

\begin{proof}
\olref[ord-arithmetic][using-addition]{ordinfinitycharacter},
\olref[cardinals][cardsasords]{lem:CardinalsBehaveRight} और
\olref{naturalsarecardinals} से।
\end{proof}

इससे उन कुछ धारणाओं की निम्न \emph{परिभाषा} उचित ठहरती है जिनका उपयोग हमने
\olref[sfr][][]{part} में कुछ अनौपचारिक ढंग से किया था:

\begin{defn}\ollabel{defnfinite}
हम कहते हैं कि $A$ \emph{परिमित} है तभी जब $\card{A}$ कोई प्राकृतिक
संख्या है, अर्थात् $\card{A}\in\omega$। अन्यथा हम कहते हैं कि $A$
\emph{अनंत} है।
\end{defn}
\noindent 
किंतु ध्यान दें कि यह परिभाषा $\ZFC$ की पृष्ठभूमि में प्रस्तुत की गई है।
आखिर प्रत्येक समुच्चय की कार्डिनलता की प्रत्याभूति के लिए हमें सुव्यवस्था
की आवश्यकता थी। और वास्तव में सुव्यवस्था के बिना कोई ऐसा समुच्चय हो सकता
है जो न परिमित हो, न डेडेकिंड-अनंत। हम \olref[choice][]{chap} में इस प्रकार
के प्रश्न पर लौटेंगे। अभी हम सुव्यवस्था पर निर्भर रहना जारी रखते हैं।

अब परिमित कार्डिनलों से अनंत कार्डिनलों की ओर बढ़ें। यहाँ दो प्राथमिक तथ्य
हैं:

\begin{cor}\ollabel{omegaisacardinal}
$\omega$ लघुतम अनंत कार्डिनल है।
\end{cor}

\begin{proof}
$\omega$ एक कार्डिनल है, क्योंकि $\omega$ डेडेकिंड-अनंत है और यदि किसी
$n\in\omega$ के लिए $\cardeq{\omega}{n}$ होता, तो $n$ डेडेकिंड-अनंत होता,
जो \olref[z][infinity-again]{naturalnumbersarentinfinite} के विरुद्ध है।
अब परिभाषा से $\omega$ लघुतम अनंत कार्डिनल है।
\end{proof}

\begin{cor}
प्रत्येक अनंत कार्डिनल एक सीमा क्रमसूचक है।
\end{cor}

\begin{proof}
$\alpha$ को एक अनंत उत्तरवर्ती क्रमसूचक मानें, अतः किसी $\beta$ के लिए
$\alpha=\beta\ordplus 1$। \olref{finitecardisoequal} से $\beta$ भी अनंत
है, अतः \olref[ord-arithmetic][using-addition]{ordinfinitycharacter} से
$\cardeq{\beta}{\beta\ordplus 1}$। अब
\olref[cardinals][cardsasords]{lem:CardinalsBehaveRight} से
$\card{\beta}=\card{\beta\ordplus 1}=\card{\alpha}$, इसलिए
$\alpha\neq\card{\alpha}$।
\end{proof}

\olref[sfr][siz][enm-alt]{defn:enumerable} में ही हमने बताया था कि हम
!!{enumerable} और !!{nonenumerable} अनंत समुच्चयों में अंतर कर सकते हैं।
वह परिभाषा स्वाभाविक रूप से निम्न परिणाम देती है:

\begin{prop}
$A$ !!{enumerable} है तभी जब $\card{A}\leq\omega$, और $A$
!!{nonenumerable} है तभी जब $\omega<\card{A}$।
\end{prop}

\begin{proof}
त्रिविभाजन से दोनों दावे तुल्य हैं, अतः यह सिद्ध करना पर्याप्त है कि $A$
!!{enumerable} है तभी जब $\card{A}\leq\omega$। \emph{दाएँ-से-बाएँ} दिशा
के लिए: यदि $\card{A}\leq\omega$, तो
\olref[cardinals][cardsasords]{lem:CardinalsBehaveRight} और
\olref{omegaisacardinal} से $\cardle{A}{\omega}$। \emph{बाएँ-से-दाएँ}
दिशा के लिए: मान लें $A$ !!{enumerable} है; तब
\olref[sfr][siz][enm-alt]{defn:enumerable} से तीन संभव स्थितियाँ हैं:
\begin{enumerate}
	\item यदि $A=\emptyset$, तो \olref{naturalsarecardinals} और
	\olref[cardinals][cardsasords]{lem:CardinalsBehaveRight} से
	$\card{A}=0\in\omega$।
	\item यदि $\cardeq{n}{A}$, तो \olref{naturalsarecardinals} और
	\olref[cardinals][cardsasords]{lem:CardinalsBehaveRight} से
	$\card{A}=n\in\omega$।
	\item यदि $\cardeq{\omega}{A}$, तो \olref{omegaisacardinal} से
	$\card{A}=\omega$।
\end{enumerate}
अतः सभी स्थितियों में $\card{A}\leq\omega$।
\end{proof}
\noindent
वास्तव में $\omega$ का स्थान विशेष है। यद्यपि अनेक गणनीय क्रमसूचक हैं:

\begin{cor}
$\omega$ एकमात्र !!{enumerable} अनंत कार्डिनल है।
\end{cor}

\begin{proof}
$\cardfont{a}$ को कोई !!a{enumerable} अनंत कार्डिनल मानें। चूँकि
$\cardfont{a}$ अनंत है, $\omega\leq\cardfont{a}$। चूँकि $\cardfont{a}$
!!a{enumerable} कार्डिनल है,
$\cardfont{a}=\card{\cardfont{a}}\leq\omega$। अतः त्रिविभाजन से
$\cardfont{a}=\omega$। \end{proof}

निस्संदेह अनंततः अनेक कार्डिनल हैं। इसलिए हम पूछ सकते हैं:
\emph{कितने कार्डिनल हैं?} निम्न परिणाम दिखाते हैं कि शायद हमें इस प्रश्न
पर पुनर्विचार करना चाहिए।

\begin{prop}\ollabel{unioncardinalscardinal}
यदि $X$ का प्रत्येक सदस्य कार्डिनल है, तो $\bigcup X$ एक कार्डिनल है।
\end{prop}

\begin{proof}
यह जाँचना सरल है कि $\bigcup X$ एक क्रमसूचक है। मान लें $\alpha\in
\bigcup X$ एक क्रमसूचक है; तब किसी कार्डिनल $\cardfont{b}$ के लिए
$\alpha\in\cardfont{b}\in X$। चूँकि $\cardfont{b}$ एक कार्डिनल है,
$\cardless{\alpha}{\cardfont{b}}$। चूँकि
$\cardfont{b}\subseteq\bigcup X$, हमारे पास
$\cardle{\cardfont{b}}{\bigcup X}$ है, और इसलिए
$\cardneq{\alpha}{\bigcup X}$। सामान्यीकरण से $\bigcup X$ एक कार्डिनल है।
\end{proof} 

\begin{thm}\ollabel{lem:NoLargestCardinal}
कोई महत्तम कार्डिनल नहीं है।
\end{thm}

\begin{proof}
किसी भी कार्डिनल $\cardfont{a}$ के लिए कैंटर का प्रमेय
(\olref[sfr][siz][car]{thm:cantor}) और
\olref[cardinals][cardsasords]{lem:CardinalsExist} से निष्पन्न होता है कि
$\cardfont{a}<\card{\Pow{\cardfont{a}}}$।
\end{proof}

\begin{thm}
सभी कार्डिनलों का समुच्चय अस्तित्व में नहीं है।
\end{thm}

\begin{proof}
विरोधाभास द्वारा मान लें $C=\Setabs{\cardfont{a}}{\cardfont{a}\text{
एक कार्डिनल है}}$। \olref{unioncardinalscardinal} से $\bigcup C$ एक
कार्डिनल है, अतः \olref{lem:NoLargestCardinal} से कोई कार्डिनल
$\cardfont{b}>\bigcup C$ है। परिभाषा से $\cardfont{b}\in C$, इसलिए
$\cardfont{b}\subseteq\bigcup{C}$, और इस प्रकार
$\cardfont{b}\leq\bigcup C$, जो विरोधाभास है।
\end{proof}

इसकी तुलना रसेल के विरोधाभास और बुराली--फ़ोर्टी दोनों से कीजिए।

\end{document}
